% ============================================================================
% sec11.tex — §11 Convective Heat Transfer, Dimensionless Groups, Dead-Air Gaps,
%             Evaporative Coupling & Top-Layer Thermal Stratification
% Target: 8 pages | Figures: fig:11:boundary-layer, fig:11:air-gap-evap
% ============================================================================
\binderchapter{11}{Convective Heat Transfer, Dimensionless Groups, Dead-Air Gaps, Evaporative Coupling \& Top-Layer Thermal Stratification}{§11 Convection \& Stratification}{sec:11}
\scopebanner{REQUIRED}[2027 Rule 5.c.iii (All Levels) \& FAQs 09/30, 10/01, 10/06 and 10/21/2026 (no stirring; probe depth; items on the cup) \textbar\ Sources: \src{S01}, \src{S02}, \src{S09}, \src{S12}, \src{S17}]

\idxtopic{convection!Newton's law of cooling}
\idxtopic{convection!natural vs forced}
\idxtopic{dimensionless numbers!Nusselt, Prandtl, Grashof, Rayleigh, Reynolds}
\idxtopic{thermal stratification!2027 cup top-layer heating}
\idxtopic{evaporative cooling!open cup mass transfer}

% ----------------------------------------------------------------------------
% L1: FORMULA & DECISION CARD
% ----------------------------------------------------------------------------
\begin{formulacard}[§11 Master Convection, Dimensionless Scaling, Air-Gap \& Evaporative Loss Decision Card \src{S02,S12}]
\begin{ftab}
$\dot{Q}_{\mathrm{conv}} = h A \left(T_s - T_\infty\right) = \dfrac{T_s - T_\infty}{R_{\mathrm{conv}}}$
  & $\dot{Q}_{\mathrm{conv}}$: convective rate $(\si{\watt})$; $h$: convection coefficient $(\si{\watt\per\square\meter\per\kelvin})$; $A$: wetted surface area $(\si{\square\meter})$; $R_{\mathrm{conv}} = 1/(hA)$ $(\si{\kelvin\per\watt})$
  & \flag{NEWTON}
  & Any solid--fluid boundary; evaluate fluid properties at $T_f = \frac{1}{2}(T_s + T_\infty)$ \idxformula{Newton's law of cooling!$\dot{Q}=hA(T_s-T_\infty)$} \\
$\mathrm{Nu}_L = \dfrac{h L}{k_{\mathrm{fluid}}} \implies h = \dfrac{k_{\mathrm{fluid}}}{L}\,\mathrm{Nu}_L$
  & $\mathrm{Nu}_L$: Nusselt number (dimensionless); $L$: characteristic length $(\si{\meter})$; $k_{\mathrm{fluid}}$: \textbf{fluid} thermal conductivity $(\si{\watt\per\meter\per\kelvin})$
  & \flag{CONV-CORE}
  & Ratio of convective heat transfer to pure fluid conduction ($\mathrm{Nu}=1$ stagnant slab; $\mathrm{Nu}_D=2$ sphere limit) \idxformula{Nusselt number!$\mathrm{Nu}_L=hL/k_{\mathrm{fluid}}$} \\
$\mathrm{Pr} = \dfrac{\nu}{\alpha} = \dfrac{\mu c_p}{k_{\mathrm{fluid}}} \approx \left(\dfrac{\delta_v}{\delta_T}\right)^{\!3}$
  & $\nu = \mu/\rho$: kinematic viscosity $(\si{\square\meter\per\second})$; $\alpha = k_f/(\rho c_p)$: thermal diffusivity $(\si{\square\meter\per\second})$; $\delta_v, \delta_T$: velocity \& thermal b.l.\ thicknesses
  & \flag{FLUID-PROP}
  & Air ($\SI{300}{\kelvin}$): $\mathrm{Pr}=0.713$; liquid water ($\SI{27}{\celsius}$): $\mathrm{Pr}=5.83$; liquid metals: $\mathrm{Pr}\sim 0.01$ \idxformula{Prandtl number!$\mathrm{Pr}=\nu/\alpha=\mu c_p/k_f$} \\
$\mathrm{Re}_L = \dfrac{u_\infty L}{\nu} = \dfrac{\rho u_\infty L}{\mu}$
  & $u_\infty$: free-stream flow speed $(\si{\meter\per\second})$; $L$: length or pipe diameter $D$ $(\si{\meter})$
  & \flag{FORCED}
  & Forced convection; laminar flat plate $\mathrm{Re}_x < \num{5e5}$, pipe $\mathrm{Re}_D < 2300$ \idxformula{Reynolds number!$\mathrm{Re}_L=u_\infty L/\nu$} \\
$\mathrm{Gr}_L = \dfrac{g \beta |T_s - T_\infty| L^3}{\nu^2}$
  & $g = \SI{9.80665}{\meter\per\square\second}$; $\beta = -\frac{1}{\rho}(\frac{\partial\rho}{\partial T})_P$ $(\si{\per\kelvin})$ [for ideal gas, $\beta = 1/T_f$ with $T_f$ in $\si{\kelvin}$!]
  & \flag{BUOYANCY}
  & Ratio of buoyancy to viscous forces; $\mathrm{Gr}_L / \mathrm{Re}_L^2 \sim 1$ indicates mixed convection \idxformula{Grashof number!$\mathrm{Gr}_L=g\beta\Delta T L^3/\nu^2$} \\
$\mathrm{Ra}_L = \mathrm{Gr}_L\,\mathrm{Pr} = \dfrac{g \beta |T_s - T_\infty| L^3}{\nu \alpha}$
  & $\mathrm{Ra}_L$: Rayleigh number (dimensionless); laminar natural convection has $\mathrm{Nu}_L \propto \mathrm{Ra}_L^{1/4}$ ($h \propto (\Delta T/L)^{1/4}$, $\dot{Q} \propto \Delta T^{5/4}$)
  & \flag{NAT-CONV}
  & Turbulent natural convection ($\mathrm{Ra}_L \gtrsim 10^9$): $\mathrm{Nu}_L \propto \mathrm{Ra}_L^{1/3}$ ($h$ independent of $L$, $\dot{Q} \propto \Delta T^{4/3}$) \idxformula{Rayleigh number!$\mathrm{Ra}_L=\mathrm{Gr}_L\mathrm{Pr}$} \\
$\mathrm{Ra}_\delta = \dfrac{g\beta\Delta T_{\mathrm{gap}}\delta^3}{\nu\alpha} < \mathrm{Ra}_c = 1708$
  & $\delta$: enclosed air-gap thickness $(\si{\meter})$; $\delta_c = \left(\dfrac{1708\,\nu\alpha}{g\beta\Delta T_{\mathrm{gap}}}\right)^{\!1/3} \approx \SIrange{10}{12}{\milli\meter}$ in room air
  & \flag{ENCLOSURE}
  & Rayleigh--B\'enard criterion: for $\delta \le \delta_c$, air stays stagnant ($\mathrm{Nu}_\delta = 1$, pure conduction $k_{\mathrm{air}}$!) \idxformula{Rayleigh-Benard threshold!$\mathrm{Ra}_c=1708$} \\
$\dot{Q}_{\mathrm{evap}} = L_v \dot{m}_{\mathrm{evap}} = L_v A_{\mathrm{surf}} g_m \left(m_{\mathrm{H_2O},s} - m_{\mathrm{H_2O},e}\right)$
  & $L_v \approx \SI{2430}{\kilo\joule\per\kilogram}$ (\SI{30}{\celsius}); $g_m = \dfrac{\rho_f D_{\mathrm{H_2O,air}}\mathrm{Nu}_{m,D}}{D}$ $(\si{\kilogram\per\square\meter\per\second})$; $m_{\mathrm{H_2O}} = \rho_{\mathrm{H_2O}}/\rho_{\mathrm{mix}}$
  & \flag{MASS-XFER}
  & Open-cup evaporative heat loss (Lienhard §11.7 Ex.~11.11 \src{S12}); suppressed $>95\%$ by BoPET film cap \idxformula{evaporative heat loss!$\dot{Q}_{\mathrm{evap}}=L_v A g_m \Delta m_{\mathrm{H_2O}}$} \\
\end{ftab}
\end{formulacard}

% ----------------------------------------------------------------------------
% L2: COMPACT EXPLANATION & BOUNDARY LAYER PHYSICS
% ----------------------------------------------------------------------------
\section{Forced vs.\ Natural Convection \& Boundary-Layer Structure}
\label{sec:11.1}

\subsection{Newton's Law of Cooling and Typical Regimes of $h$}
Convection is energy transfer between a solid surface at $T_s$ and an adjacent moving fluid at bulk temperature $T_\infty$ through the combined action of molecular thermal conduction at the stationary wall interface ($u_{\mathrm{fluid}}(y=0) = 0$, no-slip boundary condition) and macroscopic fluid advection \src{S12}. Consequently, the convective heat transfer coefficient $h$ $[\si{\watt\per\square\meter\per\kelvin}]$ in Newton's law of cooling, $\dot{Q}_{\mathrm{conv}} = h A (T_s - T_\infty)$, is \textbf{not a intrinsic thermophysical property of the fluid alone}: it depends on the surface temperature gradient inside the fluid boundary layer, $h = \frac{-k_{\mathrm{fluid}}(\partial T/\partial y)_{y=0}}{T_s - T_\infty} \sim \frac{k_{\mathrm{fluid}}}{\delta_T}$. Per Lienhard Table~1.1 (p.~21 \src{S12}), representative orders of magnitude for $h$ are:
\begin{itemize}
  \item \textbf{Natural (Free) Convection in Air / Gases:} $h \approx \SIrange{3}{20}{\watt\per\square\meter\per\kelvin}$ (around the 2027 collector box and cup in still room air: $h \approx \SIrange{4}{9}{\watt\per\square\meter\per\kelvin}$, comparable to linearized thermal radiation $h_{\mathrm{rad}} = 4\varepsilon\sigma T_m^3 \approx \SI{6}{\watt\per\square\meter\per\kelvin}$!).
  \item \textbf{Forced Convection in Air / Gases:} $h \approx \SIrange{10}{250}{\watt\per\square\meter\per\kelvin}$ (room HVAC draft or fan).
  \item \textbf{Natural Convection in Water / Liquids:} $h \approx \SIrange{50}{1000}{\watt\per\square\meter\per\kelvin}$ (bottom-heated water plume; suppressed when heated from above!).
  \item \textbf{Forced Convection in Water / Liquids:} $h \approx 100\text{--}\SI{15000}{\watt\per\square\meter\per\kelvin}$ (stirred probe in water: $h \approx 400\text{--}\SI{1200}{\watt\per\square\meter\per\kelvin}$).
  \item \textbf{Phase-Change Boiling or Condensation of Water:} $h \approx \SIrange{2500}{100000}{\watt\per\square\meter\per\kelvin}$.
\end{itemize}

\subsection{Velocity vs.\ Thermal Boundary Layers and the Role of Prandtl Number ($\mathrm{Pr}$)}
\idxtopic{boundary layer!velocity vs thermal}
As illustrated in \fref{fig:11:boundary-layer}(a), when a fluid flows along a heated surface of distance $x$ from the leading edge, viscous shear retards the fluid within a \textbf{velocity (hydrodynamic) boundary layer} of thickness $\delta_v(x)$, while heat diffuses across a \textbf{thermal boundary layer} of thickness $\delta_T(x)$ (Lienhard §6.2--6.5, pp.~292--326 \src{S12}):
\begin{equation}
  \frac{\delta_v(x)}{x} = \frac{4.92}{\sqrt{\mathrm{Re}_x}}, \quad \frac{\delta_T(x)}{\delta_v(x)} \approx \mathrm{Pr}^{-1/3}\ (\mathrm{Pr} \gtrsim 0.6), \quad \mathrm{Nu}_x = 0.332\,\mathrm{Re}_x^{1/2}\mathrm{Pr}^{1/3}, \quad \overline{\mathrm{Nu}}_L = 2\,\mathrm{Nu}_{x=L} = 0.664\,\mathrm{Re}_L^{1/2}\mathrm{Pr}^{1/3}
  \label{eq:11:blasius-pohlhausen}
\end{equation}
Notice that for laminar forced flow over a flat plate, $h_x \propto x^{-1/2}$, so the \textbf{average heat transfer coefficient $\bar{h}_L$ over length $[0, L]$ is exactly twice the local coefficient $h_x(L)$ at the trailing edge}! Moreover, the Prandtl number $\mathrm{Pr} = \nu/\alpha$ directly sets the thickness ratio $\delta_v/\delta_T \approx \mathrm{Pr}^{1/3}$:
\begin{itemize}
  \item \textbf{Gases ($\mathrm{Pr} \approx 0.71$):} $\delta_T \approx 1.12\,\delta_v$ (momentum and heat diffuse at nearly identical rates).
  \item \textbf{Liquid Water ($\mathrm{Pr} = 5.83$ at $\SI{27}{\celsius}$, $4.34$ at $\SI{40}{\celsius}$ \src{S12}) \& Oils ($\mathrm{Pr} \sim 10^2\text{--}10^4$):} $\delta_T \approx 0.56\,\delta_v$ (water) to $0.05\,\delta_v$ (oils) — heat diffuses much more slowly than momentum ($\alpha \ll \nu$), creating a thin thermal boundary layer.
  \item \textbf{Liquid Metals ($\mathrm{Hg}, \mathrm{Na}, \mathrm{Pb\text{-}Bi}$: $\mathrm{Pr} \approx 0.004\text{--}0.025$):} $\delta_T \gg \delta_v$ (free electrons conduct heat far beyond the viscous layer; $\mathrm{Nu}_x = 0.565\,\mathrm{Pe}_x^{1/2}$ where P\'eclet number $\mathrm{Pe}_x = \mathrm{Re}_x\mathrm{Pr} = u_\infty x/\alpha$).
\end{itemize}

% ---------------- FIGURE 30: fig:11:boundary-layer ---------------------------
\begin{center}
\begin{tikzpicture}[font=\scriptsize, >=Latex, scale=0.94]
  % Panel (a): Forced Convection Velocity & Thermal Boundary Layers
  \node[font=\sffamily\bfseries\scriptsize, color=navy] at (3.6, 3.85) {(a) Forced Convection Boundary Layers ($\delta_v/\delta_T \approx \mathrm{Pr}^{1/3}$)};
  \fill[gray!25] (0,0.2) rectangle (7.2,0.5);
  \draw[thick] (0,0.5) -- (7.2,0.5);
  \node[below] at (3.6, 0.22) {Heated Isothermal Flat Plate at $T_s > T_\infty$ ($x=0$ to $x=L$)};

  % Velocity boundary layer delta_v(x)
  \draw[navy, very thick] (0,0.5) .. controls (1.2,1.35) and (2.6,1.65) .. (4.2,1.80)
    .. controls (4.8,2.15) and (5.8,2.75) .. (7.1,3.05);
  \node[navy, font=\bfseries\scriptsize, above left] at (7.05, 3.02) {$\delta_v(x)$ (Velocity b.l.)};

  % Thermal boundary layer delta_T(x) for Pr > 1 (water) and Pr = 0.71 (air)
  \draw[rust, very thick, densely dashed] (0,0.5) .. controls (1.2,1.52) and (2.6,1.88) .. (4.2,2.05)
    .. controls (4.8,2.42) and (5.8,3.05) .. (7.1,3.38);
  \node[rust, font=\bfseries\scriptsize, above] at (5.3, 3.15) {$\delta_T(x)$ (Air, $\mathrm{Pr}=0.71$)};

  \draw[teal, very thick, densely dotted] (0,0.5) .. controls (1.2,1.00) and (2.6,1.18) .. (4.2,1.28)
    .. controls (4.8,1.50) and (5.8,1.85) .. (7.1,2.05);
  \node[teal, font=\bfseries\scriptsize, right] at (4.4, 1.25) {$\delta_T(x)$ (Water, $\mathrm{Pr}=5.8$)};

  % Transition line at x_c
  \draw[gray!70, thick, dashed] (4.2,0.5) -- (4.2,3.5);
  \node[align=center, fill=white, inner sep=1pt] at (2.1, 3.35) {Laminar ($\mathrm{Re}_x < \num{5e5}$)\\[-1pt]$h_x \propto x^{-1/2}$};
  \node[align=center, fill=white, inner sep=1pt] at (5.75, 0.90) {Turbulent ($\mathrm{Re}_x > \num{5e5}$)\\[-1pt]$h_x \propto x^{-1/5}$};

  % Panel (b): Bottom-Heated Unstable Plume vs. 2027 Top-Heated Stable Stratification
  \node[font=\sffamily\bfseries\scriptsize, color=navy] at (11.7, 3.85) {(b) Unstable Bottom Heating vs.\ 2027 Stable Top-Down Lamp Heating};

  % Left cup: Bottom-heated (Unstable Ra > 0)
  \draw[thick] (8.3,2.9) -- (8.55,0.5) -- (10.45,0.5) -- (10.7,2.9);
  \fill[rust!18] (8.55,0.5) -- (10.45,0.5) -- (10.63,2.2) -- (8.37,2.2) -- cycle;
  \draw[rust, ultra thick] (8.5,0.42) -- (10.5,0.42);
  \node[rust, font=\bfseries\scriptsize, below] at (9.5, 0.40) {Bottom Hot ($\rho_{\mathrm{bot}} < \rho_{\mathrm{top}}$)};
  \draw[->, rust, thick] (9.1,0.75) .. controls (8.8,1.4) and (9.0,1.9) .. (9.3,1.95) .. controls (9.5,1.8) and (9.4,1.0) .. (9.1,0.75);
  \draw[->, rust, thick] (9.9,0.75) .. controls (10.2,1.4) and (10.0,1.9) .. (9.7,1.95) .. controls (9.5,1.8) and (9.6,1.0) .. (9.9,0.75);
  \node[align=center, font=\bfseries\scriptsize, text=rust] at (9.5, 2.55) {Unstable ($\mathrm{Ra}>0$)\\[-1pt]Plumes mix entire cup};

  % Right cup: Top-heated by 2027 Heat Lamp (Stable Ra < 0, FAQs 09/30--10/06/2026 on stirring and probe depth)
  \draw[thick] (11.8,2.9) -- (12.05,0.5) -- (13.95,0.5) -- (14.2,2.9);
  % Stratified bands
  \fill[navy!12] (12.05,0.5) -- (13.95,0.5) -- (14.07,1.65) -- (11.93,1.65) -- cycle;
  \fill[rust!35] (11.93,1.65) -- (14.07,1.65) -- (14.13,2.20) -- (11.87,2.20) -- cycle;
  % BoPET glazing cap on rim
  \draw[teal, ultra thick] (11.7,2.92) -- (14.3,2.92);
  \node[teal, font=\bfseries\scriptsize, above] at (13.0, 2.93) {BoPET Cap (FAQ 10/21/2026 on items on the cup: $\dot{Q}_{\mathrm{evap}}\to 0$)};
  % Downward NIR rays
  \foreach \xx in {12.3, 12.75, 13.25, 13.7} {
    \draw[->, amber, thick] (\xx, 3.45) -- (\xx, 2.25);
  }
  \node[font=\bfseries\scriptsize, text=rust] at (13.0, 1.92) {Hot Top Layer ($T_{\mathrm{top}},\; \rho_{\mathrm{top}}$ low)};
  \node[font=\scriptsize, text=navy] at (13.0, 1.05) {Cooler Bottom ($\rho_{\mathrm{bot}}$ high)};
  \node[align=center, font=\bfseries\scriptsize, text=navy] at (13.0, 0.15) {Buoyantly Stable ($\mathrm{Ra}<0$, FAQs 09/30--10/06/2026 on stirring and probe depth)};
\end{tikzpicture}
\captionof{figure}{(a) Hydrodynamic ($\delta_v$) and thermal ($\delta_T$) boundary-layer development along a heated flat plate \src{S12}. (b) Contrast between unstable bottom heating ($\mathrm{Ra}>0$, full convective mixing) and 2027 top-down near-IR lamp heating ($\mathrm{Ra}<0$, buoyantly stable hot top water layer measured by the ES under FAQs 09/30--10/06/2026 on stirring and probe depth \src{S02}).}
\label{fig:11:boundary-layer}
\idxfig{boundary layer!velocity vs thermal and top-layer stratification}
\end{center}

% ----------------------------------------------------------------------------
% L3: VERIFIED CORRELATIONS, ENCLOSED AIR GAPS & EVAPORATIVE COUPLING
% ----------------------------------------------------------------------------
\section{Verified Natural \& Forced Convection Correlations (Lienhard Ch.\ 6--8)}
\label{sec:11.2}

\begin{table}[htbp]
\centering
\caption{Verified Forced and Natural Convection Nusselt Correlations from Lienhard \& Lienhard (6th ed., 2024, Chapters 6--8 \src{S12}). All fluid properties are evaluated at film temperature $T_f = (T_s + T_\infty)/2$ except internal pipe flows (evaluated at bulk mean $T_b$).}
\label{tab:11:correlations}
\footnotesize
\setlength{\tabcolsep}{3.2pt}
\renewcommand{\arraystretch}{1.16}
\begin{tabularx}{\linewidth}{@{}L{0.25\linewidth} L{0.10\linewidth} X L{0.25\linewidth}@{}}
\toprule
\textbf{Geometry \& Flow Regime} & \textbf{Length $L$} & \textbf{Verified Nusselt Correlation (Lienhard 6e \src{S12})} & \textbf{Validity \& Scaling} \\
\midrule
\textbf{Forced Flat Plate} (Laminar, isothermal) & $L = \text{plate length}$ & $\overline{\mathrm{Nu}}_L = 0.664\,\mathrm{Re}_L^{1/2}\mathrm{Pr}^{1/3}$ \quad (Eq.\ 6.58, p.\ 319) & $\mathrm{Re}_L < \num{5e5},\; \mathrm{Pr} \ge 0.6$; $\bar{h} \propto u_\infty^{1/2} L^{-1/2}$ \\
\textbf{Forced Flat Plate} (Turbulent, isothermal) & $L = \text{plate length}$ & $\overline{\mathrm{Nu}}_L = 0.037\,\mathrm{Re}_L^{4/5}\mathrm{Pr}^{1/3}$ \quad (Ch.\ 6, p.\ 344) & $\num{5e5} < \mathrm{Re}_L < 10^8$; $\bar{h} \propto u_\infty^{4/5} L^{-1/5}$ \\
\textbf{Internal Pipe Flow} (Laminar, fully developed) & $D = \text{tube ID}$ & $\mathrm{Nu}_D = 3.657$ (const $T_w$); $\mathrm{Nu}_D = 4.364$ (const $q''_w$) (pp.\ 371--373) & $\mathrm{Re}_D < 2300$; $h \propto k_f / D$ (independent of $u_m$!) \\
\textbf{Internal Pipe Flow} (Turbulent, Dittus--Boelter) & $D = \text{tube ID}$ & $\mathrm{Nu}_D = 0.023\,\mathrm{Re}_D^{0.8}\mathrm{Pr}^n$ ($n=0.4$ heating, $0.3$ cooling, p.\ 383) & $\mathrm{Re}_D > 10^4,\; 0.7 \le \mathrm{Pr} \le 160$; $h \propto u_m^{0.8} D^{-0.2}$ \\
\textbf{Vertical Plate / Cylinder} (Natural convection) & $L = \text{vertical height}$ & $\overline{\mathrm{Nu}}_L = 0.68 + \dfrac{0.670\,\mathrm{Ra}_L^{1/4}}{\left[1 + (0.492/\mathrm{Pr})^{9/16}\right]^{4/9}}$ (Churchill--Chu Eq.\ 8.13a, p.\ 422) & $0 < \mathrm{Ra}_L \le 10^9$ (laminar); $\bar{h} \propto (\Delta T/L)^{1/4}$, $\dot{Q} \propto \Delta T^{5/4}$ \\
\textbf{Horizontal Cylinder} (Natural convection) & $D = \text{outer diameter}$ & $\overline{\mathrm{Nu}}_D = 0.36 + \dfrac{0.518\,\mathrm{Ra}_D^{1/4}}{\left[1 + (0.559/\mathrm{Pr})^{9/16}\right]^{4/9}}$ (Churchill--Chu Eq.\ 8.28, p.\ 430) & $10^{-6} \le \mathrm{Ra}_D \le 10^9$; $\bar{h} \propto (\Delta T/D)^{1/4}$ \\
\textbf{Submerged Sphere} (Natural convection) & $D = \text{sphere diameter}$ & $\overline{\mathrm{Nu}}_D = 2 + \dfrac{0.589\,\mathrm{Ra}_D^{1/4}}{\left[1 + (0.492/\mathrm{Pr})^{9/16}\right]^{4/9}}$ (Churchill Eq.\ 8.33, p.\ 433) & $\mathrm{Ra}_D < 10^{12}$; pure conduction limit $\lim_{\mathrm{Ra}_D\to 0}\mathrm{Nu}_D = 2$ \\
\textbf{Horizontal Plate — Unstable} (Hot top OR cold bottom) & $L^* = \dfrac{A}{P}$ ($D/4$ disk) & $\overline{\mathrm{Nu}}_{L^*} = \dfrac{0.560\,\mathrm{Ra}_{L^*}^{1/4}}{\left[1 + (0.492/\mathrm{Pr})^{9/16}\right]^{4/9}}$ ($\mathrm{Ra}_{L^*} < 10^7$, Eq.\ 8.37a, p.\ 436); $\overline{\mathrm{Nu}}_{L^*} = 0.14\,\mathrm{Ra}_{L^*}^{1/3}$ ($\mathrm{Ra}_{L^*} \ge 10^7$, Eq.\ 8.35) & Buoyant plumes rise off hot surface; for $\mathrm{Ra}^{1/3}$, $\bar{h}$ is independent of $L^*$ ($\dot{Q} \propto \Delta T^{4/3}$) \\
\textbf{Horizontal Plate — Stable} (Hot bottom OR cold top) & $D = \text{disk diameter}$ & $\overline{\mathrm{Nu}}_D = 0.82\,\mathrm{Ra}_D^{1/5}\mathrm{Pr}^{0.034}$ (Kadambi--Drake Eq.\ 8.39, p.\ 437); 2D plate: $\overline{\mathrm{Nu}}_L = 0.58\,\mathrm{Ra}_L^{1/5}$ (Eq.\ 8.38) & Stable stratification pins fluid against surface ($\bar{h}$ is $3\text{--}5\times$ smaller than unstable case!) \\
\bottomrule
\end{tabularx}
\end{table}

\subsection{Internal Pipe \& Non-Circular Duct Convection, Hydraulic Diameter ($D_h$) and LMTD}
\idxtopic{internal pipe flow!constant wall temperature vs constant heat flux}
\idxtopic{log-mean temperature difference (LMTD)}
\idxtopic{hydraulic diameter!$D_h=4A_c/P_w$}
For internal flow of mass flow rate $\dot{m} = \rho u_m A_c$ $[\si{\kilogram\per\second}]$ through a duct of cross-sectional area $A_c$, wetted perimeter $P_w$, and heated length $L$ ($A_s = P_w L$, Lienhard Ch.~3 \& Ch.~7, pp.~117--134, 366--399 \src{S12}):
\begin{enumerate}
  \item \textbf{Hydraulic Diameter ($D_h$):} Used as the characteristic length in $\mathrm{Re}_{D_h} = u_m D_h / \nu$ and $\mathrm{Nu}_{D_h} = h D_h / k_f$ for any non-circular cross-section:
  \begin{equation}
    D_h \equiv \frac{4 A_c}{P_w} = \begin{cases} D & \text{(circular tube of diameter } D\text{)} \\ a & \text{(square duct of side } a\text{: } 4a^2/(4a) = a\text{)} \\ \dfrac{2ab}{a + b} & \text{(rectangular duct } a\times b\text{; wide parallel plates } b \gg a \implies D_h = 2a\text{)} \\ D_o - D_i = 2(r_o - r_i) & \text{(concentric circular annulus of radii } r_i, r_o\text{)} \end{cases}
    \label{eq:11:hydraulic-diameter}
  \end{equation}
  In fully developed laminar duct flow ($\mathrm{Re}_{D_h} < 2300$), $\mathrm{Nu}_{D_h}$ is a pure geometric constant independent of $\mathrm{Re}_{D_h}$ and $\mathrm{Pr}$ (Lienhard Table~7.1, p.~397 \src{S12}): circular tube $\mathrm{Nu}_D = 3.657$ (const $T_w$) or $4.364$ (const $q''_w$); square duct $\mathrm{Nu}_{D_h} = 2.976$ (const $T_w$) or $3.608$ (const $q''_w$); wide parallel plates $\mathrm{Nu}_{D_h} = 7.541$ (const $T_w$) or $8.235$ (const $q''_w$).
  \item \textbf{Constant Wall Heat Flux ($q''_w = \mathrm{const}$) vs.\ Constant Wall Temperature ($T_w = \mathrm{const}$) \& Log-Mean Temperature Difference ($\mathrm{LMTD}$):}
  \begin{itemize}
    \item \emph{Constant Wall Heat Flux ($q''_w = \mathrm{const}$, e.g., electrical resistance heating):} Both the bulk fluid temperature $T_b(x)$ and wall temperature $T_w(x)$ rise \textbf{linearly} at the exact same slope $dT_b/dx = (q''_w P_w)/(\dot{m} c_p)$ once thermally developed ($T_w(x) - T_b(x) = q''_w / h = \mathrm{const}$).
    \item \emph{Constant Wall Temperature ($T_w = \mathrm{const}$, e.g., condensing steam jacket or isothermal bath):} The temperature difference $\Delta T(x) = T_w - T_b(x)$ decays \textbf{exponentially} along the pipe length as the fluid approaches $T_w$:
    \begin{equation}
      \frac{\Delta T_{\mathrm{out}}}{\Delta T_{\mathrm{in}}} = \frac{T_w - T_{b,\mathrm{out}}}{T_w - T_{b,\mathrm{in}}} = \exp\!\left(-\frac{\bar{h} A_s}{\dot{m} c_p}\right) = e^{-\mathrm{NTU}}, \qquad \dot{Q} = \dot{m} c_p\left(T_{b,\mathrm{out}} - T_{b,\mathrm{in}}\right) = \bar{h} A_s \Delta T_{\mathrm{lm}}
      \label{eq:11:pipe-lmtd}
    \end{equation}
    where $\mathrm{NTU} = \frac{\bar{h}A_s}{\dot{m}c_p}$ is the Number of Transfer Units and $\Delta T_{\mathrm{lm}}$ is the \textbf{Log-Mean Temperature Difference}:
    \begin{equation}
      \boxed{\Delta T_{\mathrm{lm}} \equiv \frac{\Delta T_{\mathrm{in}} - \Delta T_{\mathrm{out}}}{\ln\!\left(\Delta T_{\mathrm{in}} / \Delta T_{\mathrm{out}}\right)} = \frac{T_{b,\mathrm{out}} - T_{b,\mathrm{in}}}{\ln\!\left(\dfrac{T_w - T_{b,\mathrm{in}}}{T_w - T_{b,\mathrm{out}}}\right)} \le \Delta T_{\mathrm{arith}} = \frac{\Delta T_{\mathrm{in}} + \Delta T_{\mathrm{out}}}{2}}
      \label{eq:11:lmtd-boxed}
    \end{equation}
  \end{itemize}
\end{enumerate}

\section{Enclosed Dead-Air Gaps ($\delta_{\mathrm{opt}} = \SIrange{10}{12}{\milli\meter}$) \& Open-Cup Evaporative Loss}
\label{sec:11.3}
\idxtopic{dead-air gap!optimal thickness 10-12 mm}
\idxtopic{Rayleigh-Benard instability!critical Rayleigh number 1708}

\subsection{Why $\SIrange{10}{12}{\milli\meter}$ Foil-Bounded Air Gaps Outperform Prohibited Styrofoam}
Consider an air layer of thickness $\delta$ trapped between two surfaces at $T_1$ and $T_2$ ($\Delta T_{\mathrm{gap}} = T_1 - T_2$). Across this gap, heat is transferred by three parallel mechanisms: gas conduction, cellular natural convection, and surface-to-surface thermal radiation:
\begin{equation}
  q''_{\mathrm{gap}}(\delta) = \mathrm{Nu}_\delta(\mathrm{Ra}_\delta)\,\frac{k_{\mathrm{air}}\Delta T_{\mathrm{gap}}}{\delta} + \frac{\sigma(T_1^4 - T_2^4)}{\frac{1}{\varepsilon_1} + \frac{1}{\varepsilon_2} - 1} \approx \underbrace{\left(\mathrm{Nu}_\delta\,k_{\mathrm{air}} + h_{\mathrm{rad}}\,\delta\right)}_{k_{\mathrm{eff}}(\delta)}\frac{\Delta T_{\mathrm{gap}}}{\delta}
  \label{eq:11:air-gap-flux}
\end{equation}
\begin{enumerate}
  \item \textbf{Rayleigh--B\'enard Convection Onset ($\mathrm{Ra}_c = 1708$):} When a horizontal fluid layer is heated from below, viscous shear prevents cellular overturning so long as $\mathrm{Ra}_\delta = \frac{g\beta\Delta T_{\mathrm{gap}}\delta^3}{\nu\alpha} < \mathrm{Ra}_c = 1708$ ($\mathrm{Nu}_\delta = 1.00$, pure stagnant conduction!). Once $\delta > \delta_c = \left(\frac{1708\,\nu\alpha}{g\beta\Delta T_{\mathrm{gap}}}\right)^{1/3}$, B\'enard convection rolls form ($\mathrm{Nu}_\delta \approx 0.069\,\mathrm{Ra}_\delta^{1/3}\mathrm{Pr}^{0.074} \propto \delta$), so making a single air gap wider than $\delta_c \approx \SI{11}{\milli\meter}$ provides \textbf{zero additional convective resistance} (plotted in \fref{fig:11:air-gap-evap}a)!
  \item \textbf{Foil Radiation Suppression:} In an unlined cardboard air gap ($\varepsilon_1 = \varepsilon_2 = 0.92$), $h_{\mathrm{rad}} \approx \SI{5.4}{\watt\per\square\meter\per\kelvin}$, which is \textbf{double} the conductive conductance $k_{\mathrm{air}}/\delta = \SI{2.62}{\watt\per\square\meter\per\kelvin}$ at $\delta = \SI{10}{\milli\meter}$! Lining the gap walls with household aluminum foil ($\varepsilon = 0.04$) drops $h_{\mathrm{rad}}$ by $41.7\times$ to $\SI{0.13}{\watt\per\square\meter\per\kelvin}$, yielding an effective conductivity $k_{\mathrm{eff}} = 0.0265 + (0.131)(0.010) = \mathbf{\SI{0.0278}{\watt\per\meter\per\kelvin}}$ — lower than banned expanded polystyrene ($\SI{0.035}{\watt\per\meter\per\kelvin}$ \src{S02,S12}) with $20\times$ lower thermal mass!
\end{enumerate}

\subsection{Open-Cup Evaporative Mass-Transfer Heat Loss (Lienhard §11.7 Example 11.11)}
Per Lienhard §11.7 (pp.~662--664 \src{S12}), an uncovered cup of warm water at surface temperature $T_s$ in ambient room air at $T_e$ and relative humidity $\mathrm{RH}$ loses latent heat rapidly because saturated water vapor above the surface ($P_{\mathrm{H_2O},s} = P_{\mathrm{sat}}(T_s)$) has lower molar mass ($M_w = \SI{18.015}{\kilogram\per\kilo\mole}$) than dry air ($M_a = \SI{28.965}{\kilogram\per\kilo\mole}$), driving upward buoyant mass transfer:
\begin{equation}
  \rho_{\mathrm{H_2O}} = \frac{P_{\mathrm{H_2O}} M_w}{R_u T}, \qquad \rho_{\mathrm{air}} = \frac{(P_{\mathrm{atm}} - P_{\mathrm{H_2O}}) M_a}{R_u T}, \qquad m_{\mathrm{H_2O}} = \frac{\rho_{\mathrm{H_2O}}}{\rho_{\mathrm{H_2O}} + \rho_{\mathrm{air}}}
  \label{eq:11:mass-fraction}
\end{equation}
\begin{equation}
  g_m = \frac{\rho_f D_{\mathrm{H_2O,air}}\,\mathrm{Nu}_{m,D}}{D_{\mathrm{surf}}}, \qquad \dot{m}_{\mathrm{evap}} = A_{\mathrm{surf}}\,g_m\left(m_{\mathrm{H_2O},s} - m_{\mathrm{H_2O},e}\right), \qquad \dot{Q}_{\mathrm{evap}} = L_v\,\dot{m}_{\mathrm{evap}}
  \label{eq:11:qevap-full}
\end{equation}
As shown in \fref{fig:11:air-gap-evap}(b), for the 2027 9 oz cup ($A_{\mathrm{surf}} = \SI{36.0}{\square\centi\meter}$) at $T_s = \SI{40}{\celsius}$--$\SI{50}{\celsius}$, \textbf{evaporation accounts for $73\text{--}83\%$ of all top-surface heat loss} ($\dot{Q}_{\mathrm{evap}} = \SIrange{2.22}{4.45}{\watt}$ vs.\ $\dot{Q}_{\mathrm{conv}} \approx \SI{0.42}{\watt}$ and $\dot{Q}_{\mathrm{rad}} \approx \SI{0.40}{\watt}$), removing $\SIrange{1330}{2670}{\joule}$ directly from the top water layer over a $\SI{10}{\minute}$ run!

% ---------------- FIGURE 31: fig:11:air-gap-evap -----------------------------
\begin{center}
\begin{tikzpicture}[font=\scriptsize]
  % Left plot: Air-gap heat flux vs thickness delta
  \begin{axis}[
    name=plot1,
    width=0.47\linewidth, height=4.5cm,
    xlabel={Enclosed Air-Gap Thickness $\delta$ $(\si{\milli\meter})$ ($\Delta T_{\mathrm{gap}} = \SI{15}{\kelvin}$)},
    ylabel={Gap Heat Flux $q''_{\mathrm{gap}}$ $(\si{\watt\per\square\meter})$},
    xmin=2, xmax=35, ymin=0, ymax=240,
    xtick={2, 6, 11, 15, 20, 25, 30, 35},
    ytick={0, 50, 100, 150, 200},
    grid=both, major grid style={gray!30},
    legend pos=north east, legend style={font=\scriptsize, fill=white, draw=navy},
    title={\textbf{(a) Dead-Air Gap Flux vs.\ Thickness $\delta$}}
  ]
    % Unlined cardboard gap (eps = 0.92 -> q_rad = 81.2 W/m2 + q_cond/conv)
    \addplot[rust, thick, densely dashed, domain=2:11, samples=30] {81.2 + 0.0265*15.0/(x*0.001)};
    \addplot[rust, thick, densely dashed, domain=11:35, samples=30, forget plot] {81.2 + 36.13};
    \addlegendentry{Cardboard ($\varepsilon=0.92$)}

    % Foil-lined air gap (eps = 0.04 -> q_rad = 1.97 W/m2 + q_cond/conv)
    \addplot[navy, ultra thick, domain=2:11, samples=40] {1.97 + 0.0265*15.0/(x*0.001)};
    \addplot[navy, ultra thick, domain=11:35, samples=30, forget plot] {1.97 + 36.13*(x/11.0)^0.05};
    \addlegendentry{Al-Foil Gap ($\varepsilon=0.04$)}

    \draw[teal, thick, dotted] (axis cs:10.94,0) -- (axis cs:10.94,160);
    \fill[navy] (axis cs:10.94,38.3) circle (2pt);
    \node[anchor=south west, text=navy, font=\bfseries\scriptsize] at (axis cs:11.2,44) {$\delta_c = \SI{10.9}{\milli\meter}$ ($\mathrm{Ra}_c=1708$)};
  \end{axis}

  % Right plot: Open-cup evaporative vs convective and radiative losses
  \begin{axis}[
    at={(plot1.outer south east)}, anchor=outer south west, xshift=0.35cm,
    width=0.47\linewidth, height=4.5cm,
    xlabel={Top Water Surface Temperature $T_s$ $(^\circ\mathrm{C})$ ($T_{\mathrm{room}}=\SI{22}{\celsius},\; 50\%\ \mathrm{RH}$)},
    ylabel={Top-Surface Heat Loss $\dot{Q}$ $(\si{\watt})$},
    xmin=22, xmax=54, ymin=0, ymax=5.5,
    xtick={22, 30, 35, 40, 45, 50, 54},
    ytick={0, 1.0, 2.0, 3.0, 4.0, 5.0},
    grid=both, major grid style={gray!30},
    legend pos=north west, legend style={font=\scriptsize, fill=white, draw=navy},
    title={\textbf{(b) Uncovered 9 oz Cup Top Losses vs.\ $T_s$}}
  ]
    % Q_evap(T_s)
    \addplot[rust, ultra thick] coordinates {(22,0.42) (26,0.72) (30,1.08) (35,1.58) (40,2.22) (45,3.08) (50,4.28) (54,5.42)};
    \addlegendentry{$\dot{Q}_{\mathrm{evap}}$ (Uncovered)}

    % Q_conv(T_s)
    \addplot[navy, thick, densely dashed, domain=22:54, samples=30] {0.0234 * (x - 22)};
    \addlegendentry{$\dot{Q}_{\mathrm{conv,top}}$}

    % Q_rad(T_s)
    \addplot[amber, thick, densely dotted, domain=22:54, samples=30] {0.0222 * (x - 22)};
    \addlegendentry{$\dot{Q}_{\mathrm{rad,top}}$ ($\varepsilon_w=0.96$)}

    % With BoPET glazing cap (FAQ 10/21/2026 on items on the cup)
    \addplot[teal, very thick, domain=22:54, samples=30] {0.011 * (x - 22)};
    \addlegendentry{$\dot{Q}_{\mathrm{tot}}$ w/ BoPET Cap}
  \end{axis}
\end{tikzpicture}
\captionof{figure}{(a) Heat flux across an enclosed air gap ($\Delta T = \SI{15}{\kelvin}$) versus gap width $\delta$, showing the Rayleigh--B\'enard onset knee at $\delta_c = \SI{10.9}{\milli\meter}$ ($\mathrm{Ra}_c = 1708$) and the $41.7\times$ radiation reduction of aluminum foil walls. (b) Uncovered 9 oz cup top-surface heat losses versus water surface temperature $T_s$ (Lienhard §11.7 \src{S12}), demonstrating why evaporation dominates and why a removable BoPET cap (FAQ 10/21/2026 on items on the cup \src{S02}) saves $\SIrange{2}{4}{\watt}$.}
\label{fig:11:air-gap-evap}
\idxfig{air gap!optimal 11 mm thickness and evaporative heat loss}
\end{center}

% ----------------------------------------------------------------------------
% L4: WORKED EXAMPLES (5 Fully Verified Examples Matching check_sec10_13.py)
% ----------------------------------------------------------------------------
\section{Worked Competition Examples (§11)}
\label{sec:11.4}

\begin{workedexample}[Natural Convection Loss from a Vertical Collector Wall (Churchill--Chu)]{Lienhard Eq.\ 8.13a \src{S12}}{ex:11.1}
A vertical side wall of a $30.0\times 30.0\times\SI{30.0}{\centi\meter}$ collector ($L = \SI{0.300}{\meter}$, width $W = \SI{0.300}{\meter}$, $A = \SI{0.0900}{\square\meter}$) has outer surface temperature $T_s = \SI{37.0}{\celsius}$ in still room air at $T_\infty = \SI{23.0}{\celsius}$. At the film temperature $T_f = \SI{30.0}{\celsius}$ ($\SI{303.15}{\kelvin}$), air has $\nu = \SI{1.60e-5}{\square\meter\per\second}$, $\alpha = \SI{2.25e-5}{\square\meter\per\second}$, $\mathrm{Pr} = 0.711$, and $k_f = \SI{0.0265}{\watt\per\meter\per\kelvin}$. Compute (a) $\mathrm{Ra}_L$, (b) $\overline{\mathrm{Nu}}_L$ and $\bar{h}$, (c) convective heat loss $\dot{Q}_{\mathrm{conv}}$, and (d) the new heat loss if $\Delta T$ doubles to $\SI{28.0}{\kelvin}$.
\tcblower
\textbf{Solution.} (a) Treating air as an ideal gas at $T_f = \SI{303.15}{\kelvin}$, $\beta = 1/T_f = \SI{3.2987e-3}{\per\kelvin}$ and $\Delta T = 37.0 - 23.0 = \SI{14.0}{\kelvin}$:
\[
  \mathrm{Ra}_L = \frac{g\beta\Delta T L^3}{\nu\alpha} = \frac{(9.80665)(\SI{3.2987e-3}{})(14.0)(0.300)^3}{(\SI{1.60e-5}{})(\SI{2.25e-5}{})} = \mathbf{\num{3.397e7}} < 10^9 \quad (\text{laminar})
\]
(b) By Churchill and Chu (\tref{tab:11:correlations}):
\[
  \overline{\mathrm{Nu}}_L = 0.68 + \frac{0.670(\num{3.3967e7})^{1/4}}{\left[1 + (0.492/0.711)^{9/16}\right]^{4/9}} = 0.68 + \frac{51.15}{1.3027} = \mathbf{39.94}, \qquad
  \bar{h} = \frac{k_f}{L}\overline{\mathrm{Nu}}_L = \frac{0.0265}{0.300}(39.94) = \mathbf{\SI{3.53}{\watt\per\square\meter\per\kelvin}}
\]
(c) Convective heat loss: $\dot{Q}_{\mathrm{conv}} = \bar{h} A \Delta T = (3.528)(0.0900)(14.0) = \mathbf{\SI{4.45}{\watt}}$.\\
(d) Doubling $\Delta T$ doubles $\mathrm{Ra}_L$ (same air properties), so $0.670\,\mathrm{Ra}_L^{1/4}$ grows by $2^{1/4}$: $\overline{\mathrm{Nu}}_L = 0.68 + (51.15)(1.1892)/1.3027 = 47.37$, $\bar{h} = \SI{4.18}{\watt\per\square\meter\per\kelvin}$, and $\dot{Q} = (4.18)(0.0900)(28.0) = \mathbf{\SI{10.5}{\watt}}$, a ratio of $2.37$, \emph{not} $2.0$. (The pure power law $\dot{Q} \propto \Delta T^{5/4}$ gives $2^{5/4} = 2.38$; the constant $0.68$ term pulls it slightly lower.)
\verifiedby{Evaluated in \texttt{check\_sec10\_13.py} (\texttt{EX-11.1-RAL}, \texttt{EX-11.1-Q}): $\mathrm{Ra}_L = \num{3.3967e7}$, $\dot{Q}_{\mathrm{conv}} = \SI{4.446}{\watt}$.}
\end{workedexample}

\begin{workedexample}[Rayleigh--B\'enard Critical Air-Gap Thickness \& Foil Baffle Conductance]{2027 Device Insulation}{ex:11.2}
An air space beneath the 9 oz cup cradle has mean air temperature $T_f = \SI{305.0}{\kelvin}$ ($\beta = 1/305.0 = \SI{3.2787e-3}{\per\kelvin},\; \nu = \SI{1.62e-5}{\square\meter\per\second},\; \alpha = \SI{2.28e-5}{\square\meter\per\second},\; k_{\mathrm{air}} = \SI{0.0265}{\watt\per\meter\per\kelvin}$) and temperature difference $\Delta T_{\mathrm{gap}} = \SI{15.0}{\kelvin}$ across one gap. Find (a) the maximum gap thickness $\delta_c$ before Rayleigh--B\'enard cellular convection begins ($\mathrm{Ra}_c = 1708$), and (b) the effective thermal conductivity $k_{\mathrm{eff}}$ of a $\delta = \SI{10.0}{\milli\meter}$ gap lined on both sides with household aluminum foil ($\varepsilon = 0.04$).
\tcblower
\textbf{Solution.} (a) Inverting $\mathrm{Ra}_\delta = \frac{g\beta\Delta T_{\mathrm{gap}}\delta_c^3}{\nu\alpha} = 1708$:
\[
  \delta_c = \left(\frac{1708\,\nu\alpha}{g\beta\Delta T_{\mathrm{gap}}}\right)^{\!1/3} = \left(\frac{1708(\SI{1.62e-5}{})(\SI{2.28e-5}{})}{(9.80665)(1/305.0)(15.0)}\right)^{\!1/3} = \mathbf{\SI{10.94}{\milli\meter}}
\]
(b) At $\delta = \SI{10.0}{\milli\meter} < \delta_c$, air is 100\% stagnant ($\mathrm{Nu}_\delta = 1.00$). With foil boundaries ($\varepsilon_1 = \varepsilon_2 = 0.04$ at $T_m = \SI{305}{\kelvin}$), the linearized radiative conductance across the gap is $h_{\mathrm{rad,foil}} = \frac{4\sigma T_m^3}{2/\varepsilon - 1} = \frac{4(\SI{5.6704e-8}{})(305.0)^3}{49.0} = \SI{0.1313}{\watt\per\square\meter\per\kelvin}$:
\[
  k_{\mathrm{eff}} = k_{\mathrm{air}} + h_{\mathrm{rad,foil}}\,\delta = 0.0265 + (0.1313)(0.0100) = \mathbf{\SI{0.02781}{\watt\per\meter\per\kelvin}}
\]
\verifiedby{Checked in \texttt{check\_sec10\_13.py} (\texttt{EX-11.2-DELTAC}, \texttt{EX-11.2-KEFF}): $\delta_c = \SI{10.936}{\milli\meter}$, $k_{\mathrm{eff}} = \SI{0.027813}{\watt\per\meter\per\kelvin}$.}
\end{workedexample}

\begin{workedexample}[Open-Cup Evaporative Mass-Transfer Heat Loss vs.\ Removable Glazing Cap]{Lienhard Ex.\ 11.11 \& FAQ 10/21/2026 on items on the cup \src{S02,S12}}{ex:11.3}
During a 2027 collector test, the top water surface in an uncovered 9 oz cup ($D_{\mathrm{surf}} = \SI{6.77}{\centi\meter},\; A_{\mathrm{surf}} = \SI{3.60e-3}{\square\meter}$) reaches $T_s = \SI{40.0}{\celsius}$ ($\SI{313.15}{\kelvin}$, saturation vapor pressure $P_{\mathrm{sat}} = \SI{7384}{\pascal}$, latent heat $L_v = \SI{2406}{\kilo\joule\per\kilogram}$ \src{S09}) in ambient room air at $T_e = \SI{22.0}{\celsius}$ ($\SI{295.15}{\kelvin}$), $\SI{50.0}{\percent}\ \mathrm{RH}$ ($P_{\mathrm{H_2O},e} = \SI{1322}{\pascal}$), and $P_{\mathrm{atm}} = \SI{101325}{\pascal}$. Given convective mass transfer conductance $g_m = \SI{6.65e-3}{\kilogram\per\square\meter\per\second}$, calculate (a) the surface and ambient water-vapor mass fractions $m_{\mathrm{H_2O},s}$ and $m_{\mathrm{H_2O},e}$, (b) the evaporation rate $\dot{m}_{\mathrm{evap}}$ and evaporative heat loss $\dot{Q}_{\mathrm{evap}}$, and (c) the temperature penalty over a $\SI{10.0}{\minute}$ ($\SI{600}{\second}$) run on the top stratified water layer ($C_{\mathrm{top}} = \SI{215}{\joule\per\kelvin}$).
\tcblower
\textbf{Solution.} (a) Using Eq.~\eref{eq:11:mass-fraction} with $R_u = \SI{8314.46}{\joule\per\kilo\mole\per\kelvin}$, $M_w = 18.0153$, and $M_a = 28.965$:
{\footnotesize\begin{align*}
  T_s = \SI{313.15}{\kelvin}:\; \rho_{\mathrm{H_2O},s} &= \frac{(7384)(18.0153)}{(8314.46)(313.15)} = \SI{0.05109}{\kilogram\per\cubic\meter},\;
  \rho_{\mathrm{air},s} = \frac{(101325 - 7384)(28.965)}{(8314.46)(313.15)} = \SI{1.0450}{\kilogram\per\cubic\meter} \Rightarrow m_{\mathrm{H_2O},s} = \mathbf{0.04661} \\
  T_e = \SI{295.15}{\kelvin}:\; \rho_{\mathrm{H_2O},e} &= \frac{(1322)(18.0153)}{(8314.46)(295.15)} = \SI{0.009704}{\kilogram\per\cubic\meter},\;
  \rho_{\mathrm{air},e} = \frac{(101325 - 1322)(28.965)}{(8314.46)(295.15)} = \SI{1.1802}{\kilogram\per\cubic\meter} \Rightarrow m_{\mathrm{H_2O},e} = \mathbf{0.008155}
\end{align*}}
(b) By Eq.~\eref{eq:11:qevap-full}, with $\Delta m_{\mathrm{H_2O}} = 0.04661 - 0.008155 = 0.03846$:
\[
  \dot{m}_{\mathrm{evap}} = (\SI{3.60e-3}{\square\meter})(\SI{6.65e-3}{\kilogram\per\square\meter\per\second})(0.03846) = \mathbf{\SI{9.206e-7}{\kilogram\per\second}} = \SI{0.921}{\milli\gram\per\second} \quad (\SI{0.552}{\gram}\text{ in }\SI{10}{\minute})
\]
\[
  \dot{Q}_{\mathrm{evap}} = L_v\,\dot{m}_{\mathrm{evap}} = (\SI{2.406e6}{\joule\per\kilogram})(\SI{9.206e-7}{\kilogram\per\second}) = \mathbf{\SI{2.215}{\watt}}
\]
(c) Over $\SI{600}{\second}$, evaporation removes $E_{\mathrm{evap}} = (\SI{2.215}{\watt})(\SI{600}{\second}) = \SI{1329}{\joule}$, depressing the top stratified layer ($C_{\mathrm{top}} = \SI{215}{\joule\per\kelvin}$) by $\Delta T_{\mathrm{loss}} = \frac{1329}{215} = \mathbf{\SI{6.18}{\celsius}}$! Resting a clear BoPET film cap over the cup mouth (FAQ 10/21/2026 on items on the cup) traps saturated vapor ($\Delta m_{\mathrm{H_2O}} \to 0$) and recovers this entire loss.
\verifiedby{Checked in \texttt{check\_sec10\_13.py} (\texttt{EX-11.3-QEVAP}): $\dot{Q}_{\mathrm{evap}} = \SI{2.2150}{\watt}$.}
\end{workedexample}

\begin{workedexample}[Top-Down Stratified Heating Advantage in Unstirred Water]{2027 Scoring Core \src{S02}}{ex:11.4}
During a $\SI{10.0}{\minute}$ heating run, a collector delivers $E_{\mathrm{net}} = \SI{5400}{\joule}$ of net retained thermal energy to the $\SI{100.0}{\milli\liter}$ water ($m_{\mathrm{tot}} = \SI{100.0}{\gram},\; c_p = \SI{4.184}{\joule\per\gram\per\kelvin},\; C_{\mathrm{bulk}} = \SI{418.4}{\joule\per\kelvin}$) starting at $T_{\mathrm{start}} = \SI{21.0}{\celsius}$. Because near-IR lamp radiation is absorbed from above ($\mathrm{Ra} < 0$, buoyantly stable), $\SI{72.4}{\percent}$ of the energy ($\SI{3910}{\joule}$) is absorbed directly in the upper $\SI{45.0}{\milli\liter}$ layer ($C_{\mathrm{top}} = \SI{188.28}{\joule\per\kelvin}$) while $\SI{27.6}{\percent}$ ($\SI{1490}{\joule}$) reaches the lower $\SI{55.0}{\milli\liter}$ layer ($C_{\mathrm{bot}} = \SI{230.12}{\joule\per\kelvin}$), and internal water conduction transfers $E_{\mathrm{cond,down}} = \SI{690}{\joule}$ from top to bottom. Compare the ES-measured top-layer temperature gain $\mathrm{TG}_{\mathrm{top}}$ (unstirred per FAQs 09/30--10/06/2026 on stirring and probe depth) against the bulk temperature gain $\mathrm{TG}_{\mathrm{bulk}}$ if a competitor accidentally swirls the cup while removing it!
\tcblower
\textbf{Solution.} Net energy retained in the top $\SI{45.0}{\milli\liter}$ layer is $E_{\mathrm{top}} = 3910 - 690 = \SI{3220}{\joule}$, while the bottom $\SI{55.0}{\milli\liter}$ layer gains $E_{\mathrm{bot}} = 1490 + 690 = \SI{2180}{\joule}$ ($E_{\mathrm{top}} + E_{\mathrm{bot}} = \SI{5400}{\joule}$):
\[
  \mathrm{TG}_{\mathrm{top}} = \frac{E_{\mathrm{top}}}{C_{\mathrm{top}}} = \frac{\SI{3220}{\joule}}{\SI{188.28}{\joule\per\kelvin}} = \mathbf{+\SI{17.10}{\celsius}} \quad (T_{f,\mathrm{top}} = \SI{38.10}{\celsius}), \qquad
  \mathrm{TG}_{\mathrm{bot}} = \frac{\SI{2180}{\joule}}{\SI{230.12}{\joule\per\kelvin}} = +\SI{9.47}{\celsius}
\]
If the competitor accidentally swirls or shakes the cup during removal (violating FAQs 09/30, 10/01 and 10/06/2026 on stirring and swirling and mixing the layers), the water equalizes to $\mathrm{TG}_{\mathrm{bulk}} = \frac{5400}{418.4} = \mathbf{+\SI{12.91}{\celsius}}$ ($T_{f,\mathrm{bulk}} = \SI{33.91}{\celsius}$) — \textbf{destroying $\SI{4.20}{\celsius}$ ($\SI{24.5}{\percent}$) of their Temperature Score $\mathrm{TS}$ AND ruining their Prediction Score $\mathrm{PS}$!}
\verifiedby{First-law enthalpy check in \texttt{check\_sec10\_13.py} (\texttt{EX-11.4-STRAT}, \texttt{EX-11.4-ECONS}): $C_{\mathrm{top}}\mathrm{TG}_{\mathrm{top}} + C_{\mathrm{bot}}\mathrm{TG}_{\mathrm{bot}} = \SI{5400.00}{\joule}$.}
\end{workedexample}

\begin{workedexample}[Mixed Convection Criterion ($\mathrm{Gr}_L/\mathrm{Re}_L^2$) Under Room HVAC Drafts]{Lienhard Eq.\ 8.45 \src{S12}}{ex:11.5}
If a tournament testing table sits beneath a room air-conditioning vent producing a gentle horizontal draft $u_\infty = \SI{0.250}{\meter\per\second}$ across the vertical wall of Worked Example~\ref{ex:11.1} ($L = \SI{0.300}{\meter},\; \nu = \SI{1.60e-5}{\square\meter\per\second},\; \mathrm{Ra}_L = \num{3.397e7},\; \mathrm{Pr} = 0.711$), evaluate $\mathrm{Gr}_L / \mathrm{Re}_L^2$ and determine whether forced convection can be neglected.
\tcblower
\textbf{Solution.} $\mathrm{Gr}_L = \frac{\mathrm{Ra}_L}{\mathrm{Pr}} = \frac{\num{3.3967e7}}{0.711} = \num{4.777e7}$, and $\mathrm{Re}_L = \frac{u_\infty L}{\nu} = \frac{(0.250)(0.300)}{\SI{1.60e-5}{}} = \mathbf{4687.5}$:
\[
  \frac{\mathrm{Gr}_L}{\mathrm{Re}_L^2} = \frac{\num{4.777e7}}{(4687.5)^2} = \frac{\num{4.777e7}}{\num{2.197e7}} = \mathbf{2.174} \sim \mathcal{O}(1)
\]
Because $\mathrm{Gr}_L/\mathrm{Re}_L^2 = 2.17$ is of order unity ($0.1 < \mathrm{Gr}_L/\mathrm{Re}_L^2 < 10$), \textbf{neither natural nor forced convection can be neglected} (mixed convection applies: $\mathrm{Nu}_{\mathrm{mixed}}^3 \approx \mathrm{Nu}_{\mathrm{forced}}^3 + \mathrm{Nu}_{\mathrm{nat}}^3$). Enclosing the collector sides within the full $\SI{30.0}{\centi\meter}$ outer shell shields the internal cup completely from room HVAC drafts!
\verifiedby{Checked in \texttt{check\_sec10\_13.py} (\texttt{EX-11.5-MIXED}): $\mathrm{Gr}_L/\mathrm{Re}_L^2 = 2.1742$.}
\end{workedexample}

\begin{workedexample}[Internal Pipe Flow Exponential Approach \& Log-Mean Temperature Difference (LMTD)]{Lienhard §7.4 \src{S12}}{ex:11.6}
Water ($c_p = \SI{4180}{\joule\per\kilogram\per\kelvin}$) flows at $\dot{m} = \SI{0.0200}{\kilogram\per\second}$ through a circular tube of inner diameter $D = \SI{15.0}{\milli\meter}$ and length $L = \SI{1.20}{\meter}$ whose inner wall is maintained at constant temperature $T_w = \SI{80.0}{\celsius}$ by condensing steam. If the water enters at $T_{b,\mathrm{in}} = \SI{20.0}{\celsius}$ and the internal convection coefficient is $\bar{h} = \SI{1850}{\watt\per\square\meter\per\kelvin}$, compute (a) the outlet water bulk temperature $T_{b,\mathrm{out}}$, (b) the total heat transfer rate $\dot{Q}$ via First-Law enthalpy balance, and (c) verify $\dot{Q} = \bar{h} A_s \Delta T_{\mathrm{lm}}$ and compare $\Delta T_{\mathrm{lm}}$ against the naive arithmetic mean $\Delta T_{\mathrm{arith}}$.
\tcblower
\textbf{Solution.} (a) Internal surface area is $A_s = \pi D L = \pi(0.0150)(1.20) = \SI{0.056549}{\square\meter}$, and stream heat-capacity rate is $\dot{m}c_p = (0.0200)(4180) = \SI{83.60}{\watt\per\kelvin}$. Number of Transfer Units: $\mathrm{NTU} = \dfrac{\bar{h}A_s}{\dot{m}c_p} = \dfrac{(1850)(0.056549)}{83.60} = \mathbf{1.2514}$. By Eq.~\eref{eq:11:pipe-lmtd}:
\[
  T_{b,\mathrm{out}} = T_w - \left(T_w - T_{b,\mathrm{in}}\right)e^{-\mathrm{NTU}} = 80.0 - (80.0 - 20.0)e^{-1.2514} = 80.0 - 60.0(0.28611) = \mathbf{\SI{62.83}{\celsius}}
\]
(b) By the First Law: $\dot{Q} = \dot{m}c_p(T_{b,\mathrm{out}} - T_{b,\mathrm{in}}) = (83.60)(62.833 - 20.00) = \mathbf{\SI{3581}{\watt}}$.\\
(c) End temperature differences are $\Delta T_{\mathrm{in}} = 80.0 - 20.0 = \SI{60.0}{\kelvin}$ and $\Delta T_{\mathrm{out}} = 80.0 - 62.833 = \SI{17.167}{\kelvin}$. By Eq.~\eref{eq:11:lmtd-boxed}:
\[
  \Delta T_{\mathrm{lm}} = \frac{60.0 - 17.167}{\ln(60.0 / 17.167)} = \frac{42.833}{1.25138} = \mathbf{\SI{34.23}{\kelvin}} \quad \left(\text{vs.\ } \Delta T_{\mathrm{arith}} = \frac{60.0 + 17.167}{2} = \SI{38.58}{\kelvin},\ +12.7\%\text{ error!}\right)
\]
Multiplying $\bar{h}A_s\Delta T_{\mathrm{lm}} = (1850)(0.056549)(34.230) = \mathbf{\SI{3581}{\watt}}$, identical to the First-Law balance!
\verifiedby{Checked in \texttt{check\_sec10\_13.py} (\texttt{EX-11.6-LMTD}): $\dot{Q} = \SI{3580.87}{\watt}$ by both methods ($\text{rel\_err} = 0.00$).}
\end{workedexample}

\begin{workedexample}[Horizontal Heated Disk: Unstable Upward-Facing vs.\ Stable Downward-Facing]{Lienhard §8.4 \src{S12}}{ex:11.7}
A flat horizontal circular plate of diameter $D = \SI{16.0}{\centi\meter}$ ($\SI{0.160}{\meter}$) is heated to $T_s = \SI{57.0}{\celsius}$ in still room air at $T_\infty = \SI{23.0}{\celsius}$ ($\Delta T = \SI{34.0}{\kelvin}$; at $T_f = \SI{40.0}{\celsius} = \SI{313.15}{\kelvin}$, air has $\beta = 1/313.15 = \SI{3.1934e-3}{\per\kelvin}$, $\nu = \SI{1.70e-5}{\square\meter\per\second}$, $\alpha = \SI{2.40e-5}{\square\meter\per\second}$, $\mathrm{Pr} = 0.708$, $k_f = \SI{0.0272}{\watt\per\meter\per\kelvin}$). Compare the natural convection coefficient $\bar{h}_{\mathrm{up}}$ on the upward-facing hot surface (unstable plumes, Eq.~8.37a with $L^* = A/P = D/4 = \SI{0.0400}{\meter}$) against $\bar{h}_{\mathrm{down}}$ on the downward-facing hot surface (stable stratification, Eq.~8.39 with $D = \SI{0.160}{\meter}$).
\tcblower
\textbf{Solution.} Diameter-based Rayleigh number is $\mathrm{Ra}_D = \frac{(9.80665)(1/313.15)(34.0)(0.160)^3}{(\SI{1.70e-5}{})(\SI{2.40e-5}{})} = \mathbf{\num{1.0689e7}}$, and $L^*$-based Rayleigh number ($L^* = D/4$) is $\mathrm{Ra}_{L^*} = \mathrm{Ra}_D / 4^3 = \mathrm{Ra}_D / 64 = \mathbf{\num{1.6702e5}}$.
\begin{enumerate}
  \item \textbf{Upward-Facing Hot Surface (Unstable, \tref{tab:11:correlations} Eq.~8.37a):}
  \[
    \overline{\mathrm{Nu}}_{L^*} = \frac{0.560(\num{1.6702e5})^{1/4}}{\left[1 + (0.492/0.708)^{9/16}\right]^{4/9}} = \frac{11.321}{1.3033} = \mathbf{8.686} \implies
    \bar{h}_{\mathrm{up}} = \frac{k_f}{L^*}\overline{\mathrm{Nu}}_{L^*} = \frac{0.0272}{0.0400}(8.686) = \mathbf{\SI{5.907}{\watt\per\square\meter\per\kelvin}}
  \]
  \item \textbf{Downward-Facing Hot Surface (Stable, \tref{tab:11:correlations} Eq.~8.39):}
  \[
    \overline{\mathrm{Nu}}_D = 0.82(\num{1.0689e7})^{1/5}(0.708)^{0.034} = \mathbf{20.630} \implies
    \bar{h}_{\mathrm{down}} = \frac{k_f}{D}\overline{\mathrm{Nu}}_D = \frac{0.0272}{0.160}(20.630) = \mathbf{\SI{3.507}{\watt\per\square\meter\per\kelvin}}
  \]
\end{enumerate}
Stable gravitational stratification reduces convective heat transfer on the downward-facing hot surface by $\mathbf{1.68\times}$!
\verifiedby{Checked in \texttt{check\_sec10\_13.py} (\texttt{EX-11.7-H-UP}, \texttt{EX-11.7-H-DOWN}): $\bar{h}_{\mathrm{up}} = \SI{5.9067}{\watt\per\square\meter\per\kelvin}$, $\bar{h}_{\mathrm{down}} = \SI{3.5071}{\watt\per\square\meter\per\kelvin}$.}
\end{workedexample}

% ----------------------------------------------------------------------------
% L5: TRAPS & DIAGNOSTICS + L6: 2027 DEVICE APPLICATION
% ----------------------------------------------------------------------------
\section{Convection Traps, Diagnostics \& 2027 Stratification Playbook}
\label{sec:11.5}

\begin{trapbox}[Top Convection \& Mass-Transfer Traps on Science Olympiad Exams]
\idxerr{convection!confusing Nusselt Nu with Biot Bi conductivity}
\idxerr{convection!using Celsius in ideal-gas beta = 1/T}
\idxerr{convection!assuming natural convection Q is linear in DeltaT}
\par\smallskip\noindent
\begin{errtab}
\textbf{1. Using solid conductivity $k_{\mathrm{solid}}$ inside $\mathrm{Nu} = hL/k_{\mathrm{fluid}}$ (or $k_{\mathrm{fluid}}$ inside $\mathrm{Bi} = hL_c/k_{\mathrm{solid}}$)}
  & Mixes up fluid boundary-layer convection ($\mathrm{Nu}$) with internal solid temperature uniformity ($\mathrm{Bi}$) (Lienhard p.~23 \src{S12}).
  & Memory lock: $\mathrm{Nu}$ = \textbf{N}ewton fluid ($k_{\mathrm{fluid}}$); $\mathrm{Bi}$ = \textbf{B}ody internal ($k_{\mathrm{body}}$). \\
\textbf{2. Plugging Celsius $T_f$ into ideal-gas expansion coefficient $\beta = 1/T_f$}
  & Overestimates $\beta$ and $\mathrm{Ra}_L$ by $\sim 10\times$ at $\SI{25}{\celsius}$!
  & Always convert $T_f$ to absolute Kelvin: $\beta_{\mathrm{gas}} = \frac{1}{T_f(\si{\kelvin})}$. \\
\textbf{3. Assuming natural convection heat loss $\dot{Q}$ scales linearly with $\Delta T$}
  & Because $h_{\mathrm{nat}} \propto \Delta T^{1/4}$ (laminar) or $\Delta T^{1/3}$ (turbulent), $\dot{Q} = hA\Delta T$ scales as $\Delta T^{5/4}$ or $\Delta T^{4/3}$!
  & If $\Delta T$ doubles in laminar natural convection, $\dot{Q}$ multiplies by $2^{5/4} = 2.378\times$. \\
\textbf{4. Making a single insulation air gap wider than $\SI{15}{\milli\meter}$}
  & Triggers Rayleigh--B\'enard cellular convection ($\mathrm{Ra}_\delta > 1708$), circulating heat across the cavity.
  & Subdivide wide cavities with aluminum-foil baffles into $\SIrange{10}{12}{\milli\meter}$ stagnant cells. \\
\end{errtab}
\end{trapbox}

\begin{rulebox}[FAQs 09/30, 10/01, 10/06 and 10/21/2026: No Stirring, Top-Layer Probe \& Items on the Cup \src{S02,S05}]
\idxdev{FAQs on stirring and probe depth!no stirring or swirling before measurement}
\idxdev{FAQ 10/21/2026 (items on the cup)!removable non-attached glazing film}
\begin{itemize}
  \item \textbf{FAQs 09/30 and 10/06/2026 (no stirring or swirling; 10/01/2026 says the same):} \emph{``No, the water will not be stirred or swirled before taking the temperature... No, the water must not be swirled, shaken or stirred before taking the temperature.''}
  \item \textbf{FAQ 10/06/2026 (probe depth):} \emph{``This will be up to the local Event Supervisor to determine the temperature measuring protocol \ldots\ It is recommended that the ES should measure the temperature of the top layer of water near the center of the cup. If the thermometer has a probe, the ES should insert it far enough into the water to take an accurate measurement, but no farther.''}
  \item \textbf{FAQ 10/21/2026 (items on the cup) \& Official ES Guide:} Teams may rest a transparent vapor barrier / glazing cover (such as a $\SI{25}{\micro\meter}$ clear BoPET disk or LDPE wrap mounted on an upper gravity ring) on or around the cup during Water Setup Time, provided \textbf{zero adhesives, tapes, rubber bands, or clips attach it to the cup}, cup shape is unaltered, and the cover is completely removed before final temperature measurement.
\end{itemize}
\end{rulebox}

% ----------------------------------------------------------------------------
% RAPID COMPETITION PROBLEM & SOLUTION BANK (§11)
% ----------------------------------------------------------------------------
\section{Rapid Competition Problem \& Solution Bank (§11)}
\label{sec:11.6}

\begin{problem}{P11.1}{Boundary-Layer Scaling Along a Laminar vs.\ Turbulent Flat Plate}{11.1}{5.c.iii}{State / National}{60 s}{Power-Law Ratio}{assuming $\bar{h}_L = h_x(L)$ instead of $2 h_x(L)$}
Air ($\mathrm{Pr} = 0.71$) flows at free-stream speed $u_\infty = \SI{4.00}{\meter\per\second}$ over an isothermal flat plate of length $L = \SI{0.400}{\meter}$ ($\nu = \SI{1.60e-5}{\square\meter\per\second}$, $k_f = \SI{0.0265}{\watt\per\meter\per\kelvin}$). (a) Verify that the boundary layer is laminar over the entire plate and compute both the trailing-edge local convection coefficient $h_x(L)$ and the plate-averaged coefficient $\bar{h}_L$. (b) If the plate length is quadrupled to $4L$ while remaining laminar, by what factor does the total heat transfer rate $\dot{Q}$ change?
\end{problem}
\begin{solution}{P11.1}
(a) At $x = L = \SI{0.400}{\meter}$, $\mathrm{Re}_L = \dfrac{u_\infty L}{\nu} = \dfrac{(4.00)(0.400)}{\SI{1.60e-5}{}} = \mathbf{\num{1.00e5} < \num{5e5}}$ (100\% laminar). By Eq.~\eref{eq:11:blasius-pohlhausen}:
\[
  \mathrm{Nu}_{x=L} = 0.332(\num{1.00e5})^{1/2}(0.71)^{1/3} = 0.332(316.23)(0.8921) = \mathbf{93.66} \implies
  h_x(L) = \frac{0.0265}{0.400}(93.66) = \mathbf{\SI{6.205}{\watt\per\square\meter\per\kelvin}}
\]
The plate-averaged coefficient is exactly twice the trailing-edge value: $\bar{h}_L = 2 h_x(L) = \mathbf{\SI{12.41}{\watt\per\square\meter\per\kelvin}}$ ($\overline{\mathrm{Nu}}_L = 187.3$).\\
(b) Since $\bar{h}_L \propto L^{-1/2}$ while surface area $A = W L \propto L^1$, total heat flow scales as $\dot{Q} = \bar{h}_L A \Delta T \propto L^{1/2}$: quadrupling $L$ multiplies total heat transfer by $4^{1/2} = \mathbf{2.00\times}$ (the upstream half $[0, 2L]$ transfers $\frac{1}{\sqrt{2}} = 70.7\%$ of the heat!).
\end{solution}

\begin{problem}{P11.2}{Microgravity / Small-Sphere Conduction Limit ($\mathrm{Nu}_D \to 2$)}{11.2}{5.c.iii}{National / MIT}{60 s}{Asymptotic Limit}{using cylinder correlation with lead constant 0.36 instead of 2}
A spherical water droplet of diameter $D = \SI{0.500}{\milli\meter}$ is suspended in still air ($k_{\mathrm{air}} = \SI{0.0262}{\watt\per\meter\per\kelvin}$) under microgravity conditions ($g \to 0 \implies \mathrm{Ra}_D \to 0$, Lienhard §8.4 Eq.~8.31, p.~433 \src{S12}). What is the exact asymptotic Nusselt number $\mathrm{Nu}_D$ and convective/conductive heat transfer coefficient $h$ of the droplet?
\end{problem}
\begin{solution}{P11.2}
As $\mathrm{Ra}_D \to 0$, buoyant motion vanishes and steady radial heat transfer into the infinite surrounding air medium occurs by pure 3-D spherical conduction ($R_{\mathrm{sph},\infty} = \lim_{r_o\to\infty}\frac{r_o - r_i}{4\pi k_{\mathrm{air}}r_i r_o} = \frac{1}{4\pi k_{\mathrm{air}}r_i} = \frac{1}{2\pi k_{\mathrm{air}}D}$). Equating $\dot{Q} = h(\pi D^2)\Delta T = 2\pi k_{\mathrm{air}}D\Delta T$ gives the exact lower bound for any sphere (\tref{tab:11:correlations}):
\[
  \lim_{\mathrm{Ra}_D \to 0} \mathrm{Nu}_D = \mathbf{2.000}, \qquad h = \frac{2 k_{\mathrm{air}}}{D} = \frac{2(\SI{0.0262}{\watt\per\meter\per\kelvin})}{\SI{0.000500}{\meter}} = \mathbf{\SI{104.8}{\watt\per\square\meter\per\kelvin}}
\]
(Notice how small diameter $D = \SI{0.5}{\milli\meter}$ creates a huge heat transfer coefficient $h = \SI{104.8}{\watt\per\square\meter\per\kelvin}$ even with zero air flow!).
\end{solution}

\begin{problem}{P11.3}{Dittus--Boelter Turbulent Pipe Flow Velocity \& Diameter Scaling}{11.2}{5.c.iii}{State / National}{60 s}{Exponents}{forgetting that mass flow $\dot{m}$ at fixed $D$ scales $u_m \propto \dot{m}$}
Water flows turbulently ($\mathrm{Re}_D = \num{4.0e4}$) through a heated circular tube of diameter $D$, giving internal heat transfer coefficient $h_1 = \SI{2400}{\watt\per\square\meter\per\kelvin}$. Using the Dittus--Boelter correlation $\mathrm{Nu}_D = 0.023\,\mathrm{Re}_D^{0.8}\mathrm{Pr}^{0.4}$, find the new heat transfer coefficient $h_2$ if the same mass flow rate $\dot{m}$ is forced through a tube of half the diameter ($D_2 = D_1 / 2$).
\end{problem}
\begin{solution}{P11.3}
At constant mass flow rate $\dot{m} = \rho u_m (\frac{\pi}{4}D^2)$, mean flow speed scales as $u_m \propto D^{-2}$, so $\mathrm{Re}_D = \frac{\rho u_m D}{\mu} = \frac{4\dot{m}}{\pi\mu D} \propto D^{-1}$. Therefore, $h = \frac{k_f}{D}\mathrm{Nu}_D \propto D^{-1}\left(\mathrm{Re}_D^{0.8}\right) \propto D^{-1}\left(D^{-0.8}\right) = \mathbf{D^{-1.8}}$! Halving the diameter at fixed mass flow rate multiplies $h$ by $2^{1.8} = \mathbf{3.482\times}$:
\[
  h_2 = 2^{1.8}(\SI{2400}{\watt\per\square\meter\per\kelvin}) = (3.4822)(2400) = \mathbf{\SI{8357}{\watt\per\square\meter\per\kelvin}}
\]
\end{solution}

\begin{problem}{P11.4}{Sling Psychrometer Wet-Bulb Adiabatic Evaporative Equilibrium}{11.3}{5.c.ii, 5.c.iii}{National / Cornell E1}{90 s}{Coupled Heat-Mass Balance}{forgetting Lewis/Chilton-Colburn factor or latent heat}
In a sling psychrometer (Lienhard §11.7 Fig.~11.16, p.~665 \src{S12}), a wet-wick thermometer bulb reaches steady wet-bulb temperature $T_{\mathrm{wb}} = \SI{18.0}{\celsius}$ when swung in dry desert air where the convective heat transfer coefficient is $h = \SI{45.0}{\watt\per\square\meter\per\kelvin}$ and the mass-transfer evaporative cooling flux at the wick surface is $\dot{m}''_{\mathrm{evap}} = \SI{0.370}{\gram\per\square\meter\per\second}$ ($L_v(\SI{18}{\celsius}) = \SI{2458}{\kilo\joule\per\kilogram}$, neglect radiation). Calculate the dry-bulb air temperature $T_{\mathrm{db}}$.
\end{problem}
\begin{solution}{P11.4}
At steady wet-bulb equilibrium, the rate of convective sensible heat gain from the warmer dry-bulb air ($q''_{\mathrm{conv}} = h(T_{\mathrm{db}} - T_{\mathrm{wb}})$) exactly balances the rate of evaporative latent heat removal ($q''_{\mathrm{evap}} = L_v \dot{m}''_{\mathrm{evap}}$):
\[
  h\left(T_{\mathrm{db}} - T_{\mathrm{wb}}\right) = L_v\,\dot{m}''_{\mathrm{evap}} = (\SI{2.458e6}{\joule\per\kilogram})(\SI{3.70e-4}{\kilogram\per\square\meter\per\second}) = \mathbf{\SI{909.46}{\watt\per\square\meter}}
\]
\[
  T_{\mathrm{db}} = T_{\mathrm{wb}} + \frac{909.46}{h} = 18.0 + \frac{909.46}{45.0} = 18.0 + 20.21 = \mathbf{\SI{38.21}{\celsius}}
\]
\end{solution}

\begin{problem}{P11.5}{Multi-Cell Foil Air-Gap Baffle Optimization Inside a 30 mm Wall}{11.3}{5.c.iii}{2027 Device / State}{75 s}{Rayleigh Scaling $\mathrm{Ra} \propto \Delta T \delta^3$}{forgetting that $N$ cells divide both $\delta$ and $\Delta T_{\mathrm{cell}}$ by $N$}
A single vertical air cavity of thickness $W = \SI{30.0}{\milli\meter}$ across a temperature drop $\Delta T_{\mathrm{tot}} = \SI{30.0}{\kelvin}$ has Rayleigh number $\mathrm{Ra}_W = \num{7.20e4}$. If two thin parallel aluminum-foil sheets are inserted to subdivide the $\SI{30.0}{\milli\meter}$ space into $N = 3$ equal $\SI{10.0}{\milli\meter}$ series air cells, by what exact factor is the Rayleigh number $\mathrm{Ra}_\delta$ inside each individual cell reduced compared to $\mathrm{Ra}_W$?
\end{problem}
\begin{solution}{P11.5}
Subdividing a gap of width $W$ into $N$ equal series cells divides the cell thickness by $N$ ($\delta = W/N$) \textbf{AND} divides the temperature drop across each cell by $N$ ($\Delta T_{\mathrm{cell}} = \Delta T_{\mathrm{tot}}/N$)! Because $\mathrm{Ra}_\delta \propto \Delta T_{\mathrm{cell}}\,\delta^3$, the Rayleigh number in each cell drops as $\mathbf{N^{-4}}$:
\[
  \frac{\mathrm{Ra}_\delta}{\mathrm{Ra}_W} = \left(\frac{\Delta T_{\mathrm{tot}}/N}{\Delta T_{\mathrm{tot}}}\right)\!\left(\frac{W/N}{W}\right)^{\!3} = \frac{1}{N^4} = \frac{1}{3^4} = \mathbf{\frac{1}{81}} \implies \mathrm{Ra}_\delta = \frac{\num{7.20e4}}{81} = \mathbf{888.9 < \mathrm{Ra}_c = 1708}
\]
Two foil sheets reduce $\mathrm{Ra}$ by $\mathbf{81\times}$, completely quenching cellular natural convection!
\end{solution}

\par\smallskip\noindent
\begin{minipage}{\linewidth}
\captionof{table}{Thermophysical Properties of Dry Air ($\SI{1}{\text{atm}}$) and Saturated Liquid Water from $\SI{0}{\celsius}$ to $\SI{100}{\celsius}$, Including Pre-Computed Rayleigh Modulus $g\beta/(\nu\alpha)$ for Instant $\mathrm{Ra}_L = [g\beta/(\nu\alpha)]\Delta T L^3$ Lookup (Lienhard Tables~A.3 \& A.6, pp.~746, 756 \src{S11,S12}).}
\label{tab:11:air-water-props}
\idxtopic{thermophysical properties!dry air 0 to 100 C}
\idxtopic{thermophysical properties!liquid water 0 to 100 C}
\idxtopic{Rayleigh modulus!pre-computed g*beta/(nu*alpha)}
\footnotesize
\setlength{\tabcolsep}{4pt}
\renewcommand{\arraystretch}{1.12}
\begin{tabularx}{\linewidth}{@{}r r r r r r r r X@{}}
\toprule
$T$ $[\si{\celsius}]$ & $T$ $[\si{\kelvin}]$ & $\rho$ $[\si{\kilogram\per\cubic\meter}]$ & $c_p$ $[\si{\joule\per\kilogram\per\kelvin}]$ & $k_f$ $[\si{\watt\per\meter\per\kelvin}]$ & $\nu$ $[\si{\square\meter\per\second}]$ & $\alpha$ $[\si{\square\meter\per\second}]$ & $\mathrm{Pr}$ $[1]$ & $g\beta/(\nu\alpha)$ $[\si{\per\kelvin\per\cubic\meter}]$ \\
\midrule
\multicolumn{9}{l}{\textbf{Dry Air at $\SI{1.00}{\text{atm}}$ ($\SI{101.325}{\kilo\pascal}$)} --- Ideal-gas thermal expansion $\beta = 1/T(\si{\kelvin})$} \\
$0$   & $273.15$ & $1.292$ & $1006$ & $0.0242$ & $\num{1.33e-5}$ & $\num{1.86e-5}$ & $0.715$ & $\num{1.45e8}$ \\
$20$  & $293.15$ & $1.204$ & $1006$ & $0.0257$ & $\num{1.51e-5}$ & $\num{2.12e-5}$ & $0.713$ & $\num{1.04e8}$ \\
$27$  & $300.00$ & $1.177$ & $1006$ & $0.0262$ & $\num{1.57e-5}$ & $\num{2.21e-5}$ & $0.712$ & $\num{9.42e7}$ \\
$40$  & $313.15$ & $1.127$ & $1007$ & $0.0272$ & $\num{1.70e-5}$ & $\num{2.40e-5}$ & $0.710$ & $\num{7.68e7}$ \\
$60$  & $333.15$ & $1.059$ & $1008$ & $0.0287$ & $\num{1.90e-5}$ & $\num{2.69e-5}$ & $0.708$ & $\num{5.76e7}$ \\
$80$  & $353.15$ & $0.999$ & $1010$ & $0.0302$ & $\num{2.10e-5}$ & $\num{2.99e-5}$ & $0.706$ & $\num{4.42e7}$ \\
$100$ & $373.15$ & $0.946$ & $1012$ & $0.0317$ & $\num{2.31e-5}$ & $\num{3.31e-5}$ & $0.704$ & $\num{3.44e7}$ \\
\midrule
\multicolumn{9}{l}{\textbf{Saturated Liquid Water ($\mathrm{H_2O}$)} --- Note $\beta < 0$ for $T < \SI{3.984}{\celsius}$ and $\beta = 0$ at $\SI{3.984}{\celsius}$} \\
$0$   & $273.15$ & $999.8$ & $4217$ & $0.561$  & $\num{1.79e-6}$ & $\num{1.33e-7}$ & $13.47$ & $\num{-2.80e9}$ \\
$20$  & $293.15$ & $998.2$ & $4182$ & $0.598$  & $\num{1.00e-6}$ & $\num{1.43e-7}$ & $7.01$  & $\num{1.42e10}$ \\
$30$  & $303.15$ & $995.6$ & $4179$ & $0.615$  & $\num{8.01e-7}$ & $\num{1.48e-7}$ & $5.42$  & $\num{2.51e10}$ \\
$40$  & $313.15$ & $992.2$ & $4179$ & $0.631$  & $\num{6.58e-7}$ & $\num{1.52e-7}$ & $4.33$  & $\num{3.77e10}$ \\
$60$  & $333.15$ & $983.2$ & $4185$ & $0.654$  & $\num{4.74e-7}$ & $\num{1.59e-7}$ & $2.98$  & $\num{6.79e10}$ \\
$80$  & $353.15$ & $971.8$ & $4197$ & $0.670$  & $\num{3.65e-7}$ & $\num{1.64e-7}$ & $2.22$  & $\num{1.05e11}$ \\
$100$ & $373.15$ & $958.4$ & $4216$ & $0.679$  & $\num{2.94e-7}$ & $\num{1.68e-7}$ & $1.75$  & $\num{1.49e11}$ \\
\bottomrule
\end{tabularx}
\end{minipage}
\par\smallskip

\begin{problem}{P11.6}{Counterflow vs.\ Parallel-Flow Heat Exchanger LMTD Comparison}{11.2}{5.c.iii}{National / Invitational}{75 s}{LMTD Eq.\ (11.5)}{using arithmetic mean $\Delta T$ or parallel-flow endpoint subtraction in counterflow}
In a double-pipe heat exchanger, hot oil enters at $T_{h,\mathrm{in}} = \SI{90.0}{\celsius}$ and exits at $T_{h,\mathrm{out}} = \SI{50.0}{\celsius}$ while transferring $\dot{Q} = \SI{12.00}{\kilo\watt}$ to cooling water that enters at $T_{c,\mathrm{in}} = \SI{20.0}{\celsius}$ and exits at $T_{c,\mathrm{out}} = \SI{45.0}{\celsius}$ (overall heat transfer coefficient $U = \SI{400.0}{\watt\per\square\meter\per\kelvin}$, Lienhard Ch.~3 pp.~117--128 \src{S12}). Compute the log-mean temperature difference $\Delta T_{\mathrm{lm}}$ and required heat-exchange surface area $A_s = \dot{Q}/(U\Delta T_{\mathrm{lm}})$ for (a) parallel flow, and (b) counterflow.
\end{problem}
\begin{solution}{P11.6}
\begin{enumerate}
  \item[(a)] \textbf{Parallel Flow (both fluids enter at End A and exit at End B):} $\Delta T_A = T_{h,\mathrm{in}} - T_{c,\mathrm{in}} = 90.0 - 20.0 = \SI{70.0}{\kelvin}$, and $\Delta T_B = T_{h,\mathrm{out}} - T_{c,\mathrm{out}} = 50.0 - 45.0 = \SI{5.00}{\kelvin}$. By Eq.~\eref{eq:11:lmtd-boxed}:
  \[
    \Delta T_{\mathrm{lm,par}} = \frac{70.0 - 5.00}{\ln(70.0/5.00)} = \frac{65.00}{\ln(14.00)} = \frac{65.00}{2.63906} = \mathbf{\SI{24.63}{\kelvin}} \implies A_{s,\mathrm{par}} = \frac{\SI{12000}{\watt}}{(400.0)(24.630)} = \mathbf{\SI{1.218}{\square\meter}}
  \]
  \item[(b)] \textbf{Counterflow (hot enters End A where cold exits; hot exits End B where cold enters):} $\Delta T_A = T_{h,\mathrm{in}} - T_{c,\mathrm{out}} = 90.0 - 45.0 = \SI{45.0}{\kelvin}$, and $\Delta T_B = T_{h,\mathrm{out}} - T_{c,\mathrm{in}} = 50.0 - 20.0 = \SI{30.0}{\kelvin}$:
  \[
    \Delta T_{\mathrm{lm,cf}} = \frac{45.0 - 30.0}{\ln(45.0/30.0)} = \frac{15.00}{\ln(1.500)} = \frac{15.00}{0.40547} = \mathbf{\SI{36.99}{\kelvin}} \implies A_{s,\mathrm{cf}} = \frac{\SI{12000}{\watt}}{(400.0)(36.995)} = \mathbf{\SI{0.8109}{\square\meter}}
  \]
  Counterflow provides a $+\SI{50.2}{\percent}$ larger driving $\Delta T_{\mathrm{lm}}$ and requires $\SI{33.4}{\percent}$ less surface area (and only counterflow allows $T_{c,\mathrm{out}} > T_{h,\mathrm{out}}$!).
\end{enumerate}
\end{solution}

\begin{problem}{P11.7}{Chilton--Colburn / Reynolds Heat-Momentum Analogy ($C_f \leftrightarrow \mathrm{St}$)}{11.1}{5.c.iii}{National / MIT}{60 s}{Chilton--Colburn Analogy}{forgetting the $\mathrm{Pr}^{2/3}$ factor when $\mathrm{Pr} \ne 1$}
Air ($\rho = \SI{1.177}{\kilogram\per\cubic\meter}$, $c_p = \SI{1006}{\joule\per\kilogram\per\kelvin}$, $\mathrm{Pr} = 0.712$ at $\SI{300}{\kelvin}$, \tref{tab:11:air-water-props}) flows at $u_\infty = \SI{12.0}{\meter\per\second}$ over a flat surface where wind-tunnel drag measurements give a plate-averaged skin-friction coefficient $\bar{C}_f = \dfrac{\bar{\tau}_w}{\frac{1}{2}\rho u_\infty^2} = 0.00420$. Using the Chilton--Colburn Reynolds analogy $\overline{\mathrm{St}}\,\mathrm{Pr}^{2/3} = \dfrac{\bar{C}_f}{2}$ where the Stanton number is $\overline{\mathrm{St}} = \dfrac{\bar{h}}{\rho c_p u_\infty} = \dfrac{\overline{\mathrm{Nu}}_L}{\mathrm{Re}_L \mathrm{Pr}}$ (Lienhard §6.3 Eq.~6.48, p.~297 \src{S12}), compute the average convective heat transfer coefficient $\bar{h}$.
\end{problem}
\begin{solution}{P11.7}
By the Chilton--Colburn analogy (valid for laminar and turbulent boundary layers with $0.6 \le \mathrm{Pr} \le 60$):
\[
  \overline{\mathrm{St}} = \frac{\bar{h}}{\rho c_p u_\infty} = \frac{\bar{C}_f}{2\,\mathrm{Pr}^{2/3}} = \frac{0.00420}{2(0.712)^{2/3}} = \frac{0.002100}{0.79736} = \mathbf{0.002634}
\]
\[
  \bar{h} = \rho c_p u_\infty\,\overline{\mathrm{St}} = (\SI{1.177}{\kilogram\per\cubic\meter})(\SI{1006}{\joule\per\kilogram\per\kelvin})(\SI{12.0}{\meter\per\second})(0.0026336) = \mathbf{\SI{37.42}{\watt\per\square\meter\per\kelvin}}
\]
\end{solution}

\begin{problem}{P11.8}{Hot-Wire / Cylinder Crossflow Anemometer Speed Determination}{11.2}{5.c.iii}{State / National}{75 s}{Churchill--Bernstein / King's Law}{forgetting that $\mathrm{Nu}_D \propto \mathrm{Re}_D^{1/2}$ at moderate $\mathrm{Re}_D$}
A cylindrical platinum resistance probe of diameter $D = \SI{1.20}{\milli\meter}$ and length $\ell = \SI{25.0}{\milli\meter}$ ($A_s = \pi D\ell = \SI{9.425e-5}{\square\meter}$) is held at $T_w = \SI{74.0}{\celsius}$ in an air stream at $T_\infty = \SI{20.0}{\celsius}$ ($T_f = \SI{47.0}{\celsius} \implies \nu = \SI{1.77e-5}{\square\meter\per\second}$, $k_f = \SI{0.0277}{\watt\per\meter\per\kelvin}$, $\mathrm{Pr} = 0.709$). When electrical heating power $P = \SI{0.720}{\watt}$ is dissipated in the probe at steady state, use the simplified Hilpert crossflow correlation $\overline{\mathrm{Nu}}_D = 0.683\,\mathrm{Re}_D^{0.466}\mathrm{Pr}^{1/3}$ ($40 < \mathrm{Re}_D < 4000$, Lienhard §8.4 \src{S12}) to determine the air speed $u_\infty$.
\end{problem}
\begin{solution}{P11.8}
From the measured electrical power and temperature difference $\Delta T = 74.0 - 20.0 = \SI{54.0}{\kelvin}$:
\[
  \bar{h} = \frac{P}{A_s \Delta T} = \frac{\SI{0.720}{\watt}}{(\SI{9.4248e-5}{\square\meter})(\SI{54.0}{\kelvin})} = \mathbf{\SI{141.47}{\watt\per\square\meter\per\kelvin}} \implies
  \overline{\mathrm{Nu}}_D = \frac{\bar{h} D}{k_f} = \frac{(141.47)(0.00120)}{0.0277} = \mathbf{6.1287}
\]
Invert $\overline{\mathrm{Nu}}_D = 0.683\,\mathrm{Re}_D^{0.466}\mathrm{Pr}^{1/3}$ with $\mathrm{Pr}^{1/3} = (0.709)^{1/3} = 0.8917$:
\[
  \mathrm{Re}_D = \left(\frac{6.1287}{0.683 \times 0.8917}\right)^{\!1/0.466} = (10.0631)^{2.1459} = \mathbf{141.8} \implies
  u_\infty = \frac{\nu\,\mathrm{Re}_D}{D} = \frac{(\SI{1.77e-5}{})(141.8)}{0.00120} = \mathbf{\SI{2.09}{\meter\per\second}}
\]
\end{solution}

\alsosee{\spref{sec:01.10} (\texttt{QS-10} Convection Fast Sheet), \spref{sec:10} (Conduction Networks), \spref{sec:12} (Thermal Radiation \& View Factors), \spref{sec:13} (Transient Lumped Heating), \spref{sec:15} (2027 Collector Stratification \& Glazing Design), \spref{sec:17} (ES Probe Immersion Depth).}
